\subsection{Reading a Matrix Through Basis Vectors}

For
\[
A=\begin{pmatrix}2&1\\[1mm]1&2\end{pmatrix},
\]
the columns are
\[
Ae_1=\begin{pmatrix}2\\1\end{pmatrix},
\qquad
Ae_2=\begin{pmatrix}1\\2\end{pmatrix}.
\]
Thus the matrix is determined by where it sends the two coordinate directions.

\begin{center}
\begin{tikzpicture}[>=Latex,scale=0.9]
  \begin{scope}[xshift=-3.5cm]
    \draw[->] (-0.4,0) -- (2.8,0) node[right] {$x$};
    \draw[->] (0,-0.4) -- (0,2.8) node[above] {$y$};
    \draw[gray!35] (0,0) grid (2,2);
    \draw[blue!75!black,very thick,->] (0,0) -- (1,0) node[midway,below] {$e_1$};
    \draw[orange!85!black,very thick,->] (0,0) -- (0,1) node[midway,left] {$e_2$};
    \node at (1.1,-0.8) {input directions};
  \end{scope}

  \draw[->,thick] (-0.9,1.2) -- (0.9,1.2) node[midway,above] {$A$};

  \begin{scope}[xshift=3.5cm]
    \draw[->] (-0.4,0) -- (3.2,0) node[right] {$x$};
    \draw[->] (0,-0.4) -- (0,3.2) node[above] {$y$};
    \filldraw[fill=gray!10,draw=gray!60] (0,0)--(2,1)--(3,3)--(1,2)--cycle;
    \draw[blue!75!black,very thick,->] (0,0) -- (2,1) node[midway,below right] {$Ae_1$};
    \draw[orange!85!black,very thick,->] (0,0) -- (1,2) node[midway,left] {$Ae_2$};
    \node at (1.4,-0.8) {output directions};
  \end{scope}
\end{tikzpicture}
\end{center}

\begin{interpretationbox}[Columns describe the action]
Every vector $x=x_1e_1+x_2e_2$ is sent to
\[
Ax=x_1Ae_1+x_2Ae_2.
\]
Matrix multiplication therefore combines the images of the basis directions with the input coordinates. The row--column calculation and the geometric action are two descriptions of the same operation.
\end{interpretationbox}
